% Generated draft: requires mathematical and source review.
\documentclass[11pt]{article}
\usepackage[margin=25mm]{geometry}
\usepackage{fontspec}
\setmainfont{Latin Modern Roman}
\usepackage{amsmath,amssymb,mathtools,unicode-math}
\setmathfont{Latin Modern Math}
\usepackage{enumitem,needspace,xurl,hyperref}
\setlength{\parindent}{0pt}
\setlength{\parskip}{5pt}
\setlength{\emergencystretch}{3em}
\newcommand{\N}{\mathbb{N}}
\newcommand{\Z}{\mathbb{Z}}
\newcommand{\Q}{\mathbb{Q}}
\newcommand{\R}{\mathbb{R}}
\newcommand{\C}{\mathbb{C}}
\newcommand{\F}{\mathbb{F}}
\newcommand{\A}{\mathbb{A}}
\newcommand{\PP}{\mathbb{P}}
\newcommand{\E}{\mathbb{E}}
\newcommand{\eps}{\varepsilon}
\newcommand{\dd}{\,\mathrm{d}}
\begin{document}
\section*{UnsolvedMath ID: 700513}
\textbf{Source problem code:} catalogue entry
\subsection{Definitions and mathematical statement}
Work in the group language consisting of multiplication, inversion, and the identity element \(1\). A group quasiidentity is a universally quantified formula
\[
\forall x_1\cdots\forall x_n\;\left(\bigwedge_{i=1}^{m}u_i(x_1,\ldots,x_n)=v_i(x_1,\ldots,x_n)\;\Longrightarrow\;u(x_1,\ldots,x_n)=v(x_1,\ldots,x_n)\right),
\]
where all terms are group terms and \(m,n\) are finite nonnegative integers. An empty conjunction is interpreted as true. Let \(\mathcal F\) denote the class of all finite groups, including the trivial group, and let
\[
Q(\mathcal F)=\{\varphi:\varphi\text{ is a group quasiidentity and }G\models\varphi\text{ for every }G\in\mathcal F\}.
\]
For a set \(\Sigma\) of group quasiidentities, write \(\Sigma\models_{\mathrm{Grp}}\varphi\) when every group satisfying every member of \(\Sigma\) satisfies \(\varphi\). Under the interpretation that a basis concerns the quasiidentities valid in all finite groups, with group axioms as background, a basis is a set \(\Sigma\subseteq Q(\mathcal F)\) such that, for every group quasiidentity \(\varphi\),
\[
\varphi\in Q(\mathcal F)\quad\Longleftrightarrow\quad\Sigma\models_{\mathrm{Grp}}\varphi.
\]
Such a basis is independent if
\[
\Sigma\setminus\{\sigma\}\not\models_{\mathrm{Grp}}\sigma
\qquad\text{for every }\sigma\in\Sigma.
\]
Equivalently, for each \(\sigma\in\Sigma\), there is a group \(H_\sigma\), not required to be finite, satisfying every member of \(\Sigma\setminus\{\sigma\}\) and failing \(\sigma\).

\textbf{Question (A. K. Rumyantsev and D. M. Smirnov, 1982).} Does \(Q(\mathcal F)\) admit an independent basis \(\Sigma\)? No finiteness or countability restriction on the basis is imposed by the supplied statement.

\textbf{Formulation qualification.} The source asks whether the class of all finite groups possesses an independent basis of quasiidentities, without defining basis or independence. The formulation above uses the interpretation of a basis for its quasi-equational theory; it does not require the models of \(\Sigma\) to be precisely the finite groups. The supplied record also does not explicitly specify whether logical consequence and independence are relative to the group axioms or to all algebras in the group language. The intended convention requires confirmation.
\subsection{Short English statement}
The question asks whether the quasiidentities valid in every finite group can be generated by a set in which no member follows from the others. This draft treats the group axioms as background and records the source's unspecified conventions for human review. The registry labels the problem open; that status has not been independently verified.
\subsection{Sources}
\begin{itemize}
\item Local UnsolvedMath problem registry, permanent ID \(700513\), exact source code \texttt{catalogue entry}, title \emph{Kourovka Notebook Problem 8.64}. The supplied record labels the problem \texttt{open} and records an update date of \texttt{2026-10-16}.
\item Supplied source attribution: \emph{The Kourovka Notebook (Unsolved Problems in Group Theory)}, 8th issue (1982), problem 8.64; proposed by A. K. Rumyantsev and D. M. Smirnov.
\item The supplied background also identifies the 21st issue edition (2026), \texttt{arXiv:1401.0300v48}. The collection description names E. I. Khukhro and V. D. Mazurov as editors of Issues 1--21 (1965--2026). These bibliographic details are supplied attribution and have not been independently checked.
\end{itemize}
\end{document}
